% nguồn: Soạn theo Lập phương của một tổng. Lập phương của một hiệu – tài liệu đính kèm
\begin{bt}%[8D2H7-1]
	Khai triển các biểu thức sau:\\
	a) $(x+2)^3$;\qquad b) $(x-3)^3$.
	\loigiai{
		a) $(x+2)^3=x^3+3\cdot x^2\cdot 2+3\cdot x\cdot 2^2+2^3=x^3+6x^2+12x+8$.\\
		b) $(x-3)^3=x^3-3\cdot x^2\cdot 3+3\cdot x\cdot 3^2-3^3=x^3-9x^2+27x-27$.
	}
\end{bt}

\begin{bt}%[8D2H7-1]
	Khai triển các biểu thức sau:\\
	a) $(2x+1)^3$;\qquad b) $(3y-2)^3$.
	\loigiai{
		a) $(2x+1)^3=(2x)^3+3\cdot(2x)^2\cdot 1+3\cdot 2x\cdot 1^2+1^3=8x^3+12x^2+6x+1$.\\
		b) $(3y-2)^3=(3y)^3-3\cdot(3y)^2\cdot 2+3\cdot 3y\cdot 2^2-2^3=27y^3-54y^2+36y-8$.
	}
\end{bt}

\begin{bt}%[8D2H7-1]
	Khai triển các biểu thức sau:\\
	a) $(x+2y)^3$;\qquad b) $(2a-b)^3$.
	\loigiai{
		a) $(x+2y)^3=x^3+3\cdot x^2\cdot 2y+3\cdot x\cdot(2y)^2+(2y)^3=x^3+6x^2y+12xy^2+8y^3$.\\
		b) $(2a-b)^3=(2a)^3-3\cdot(2a)^2\cdot b+3\cdot 2a\cdot b^2-b^3=8a^3-12a^2b+6ab^2-b^3$.
	}
\end{bt}

\begin{bt}%[8D2H7-1]
	Khai triển các biểu thức sau:\\
	a) $\left(\dfrac12x-2\right)^3$;\qquad b) $\left(x^2+1\right)^3$.
	\loigiai{
		a) $\left(\dfrac12x-2\right)^3=\left(\dfrac12x\right)^3-3\cdot\left(\dfrac12x\right)^2\cdot 2+3\cdot\dfrac12x\cdot 2^2-2^3$\\
		$=\dfrac18x^3-\dfrac32x^2+6x-8$.\\
		b) $\left(x^2+1\right)^3=\left(x^2\right)^3+3\cdot\left(x^2\right)^2\cdot 1+3\cdot x^2\cdot 1^2+1^3=x^6+3x^4+3x^2+1$.
	}
\end{bt}

\begin{bt}%[8D2H7-2]
	Viết các biểu thức sau dưới dạng lập phương của một tổng hoặc một hiệu:\\
	a) $x^3+3x^2+3x+1$;\qquad b) $x^3-6x^2+12x-8$.
	\loigiai{
		a) $x^3+3x^2+3x+1=x^3+3\cdot x^2\cdot 1+3\cdot x\cdot 1^2+1^3=(x+1)^3$.\\
		b) $x^3-6x^2+12x-8=x^3-3\cdot x^2\cdot 2+3\cdot x\cdot 2^2-2^3=(x-2)^3$.
	}
\end{bt}

\begin{bt}%[8D2H7-2]
	Viết các biểu thức sau dưới dạng lập phương của một tổng hoặc một hiệu:\\
	a) $8x^3+12x^2+6x+1$;\qquad b) $27-27y+9y^2-y^3$.
	\loigiai{
		a) $8x^3+12x^2+6x+1=(2x)^3+3\cdot(2x)^2\cdot 1+3\cdot 2x\cdot 1^2+1^3=(2x+1)^3$.\\
		b) $27-27y+9y^2-y^3=3^3-3\cdot 3^2\cdot y+3\cdot 3\cdot y^2-y^3=(3-y)^3$.
	}
\end{bt}

\begin{bt}%[8D2H7-2]
	Viết các biểu thức sau dưới dạng lập phương của một tổng hoặc một hiệu:\\
	a) $x^3+9x^2y+27xy^2+27y^3$;\qquad b) $8a^3-36a^2b+54ab^2-27b^3$.
	\loigiai{
		a) $x^3+9x^2y+27xy^2+27y^3=x^3+3\cdot x^2\cdot 3y+3\cdot x\cdot(3y)^2+(3y)^3=(x+3y)^3$.\\
		b) $8a^3-36a^2b+54ab^2-27b^3=(2a)^3-3\cdot(2a)^2\cdot 3b+3\cdot 2a\cdot(3b)^2-(3b)^3=(2a-3b)^3$.
	}
\end{bt}

\begin{bt}%[8D2H7-2]
	Điền biểu thức thích hợp vào chỗ trống ($\ldots$) để được hằng đẳng thức đúng:\\
	a) $x^3+\ldots+48x+64=(x+4)^3$;\\
	b) $27a^3-27a^2+\ldots-1=(\ldots-\ldots)^3$.
	\loigiai{
		a) $(x+4)^3=x^3+3\cdot x^2\cdot 4+3\cdot x\cdot 4^2+4^3=x^3+12x^2+48x+64$.\\
		Vậy biểu thức cần điền là $12x^2$.\\
		b) $27a^3=(3a)^3$, $1=1^3$ nên ta xét $(3a-1)^3$:
		\[ (3a-1)^3=(3a)^3-3\cdot(3a)^2\cdot 1+3\cdot 3a\cdot 1^2-1^3=27a^3-27a^2+9a-1. \]
		Vậy $27a^3-27a^2+9a-1=(3a-1)^3$.
	}
\end{bt}

\begin{bt}%[8D2H7-3]
	Tính nhanh:\\
	a) $101^3$;\qquad b) $99^3$.
	\loigiai{
		a) $101^3=(100+1)^3=100^3+3\cdot 100^2\cdot 1+3\cdot 100\cdot 1^2+1^3$\\
		$=1\,000\,000+30\,000+300+1=1\,030\,301$.\\
		b) $99^3=(100-1)^3=100^3-3\cdot 100^2\cdot 1+3\cdot 100\cdot 1^2-1^3$\\
		$=1\,000\,000-30\,000+300-1=970\,299$.
	}
\end{bt}

\begin{bt}%[8D2H7-3]
	Tính giá trị của các biểu thức sau:\\
	a) $A=x^3+3x^2+3x+1$ tại $x=99$;\\
	b) $B=x^3-6x^2+12x-8$ tại $x=12$.
	\loigiai{
		a) $A=(x+1)^3$. Tại $x=99$: $A=(99+1)^3=100^3=1\,000\,000$.\\
		b) $B=(x-2)^3$. Tại $x=12$: $B=(12-2)^3=10^3=1\,000$.
	}
\end{bt}

\begin{bt}%[8D2H7-3]
	Rút gọn các biểu thức sau:\\
	a) $A=(x+1)^3-(x-1)^3$;\\
	b) $B=(x-2)^3-x^3+6x^2$.
	\loigiai{
		a) $A=\left(x^3+3x^2+3x+1\right)-\left(x^3-3x^2+3x-1\right)$\\
		$=x^3+3x^2+3x+1-x^3+3x^2-3x+1=6x^2+2$.\\
		b) $B=\left(x^3-6x^2+12x-8\right)-x^3+6x^2=12x-8$.
	}
\end{bt}

\begin{bt}%[8D2H7-3]
	Rút gọn các biểu thức sau:\\
	a) $C=(2x+1)^3-8x^3-12x^2$;\\
	b) $D=(x-y)^3+3xy(x-y)$.
	\loigiai{
		a) $C=\left(8x^3+12x^2+6x+1\right)-8x^3-12x^2=6x+1$.\\
		b) $D=\left(x^3-3x^2y+3xy^2-y^3\right)+\left(3x^2y-3xy^2\right)=x^3-y^3$.
	}
\end{bt}

\begin{bt}%[8D2V7-1]
	Một khối gỗ hình lập phương có cạnh $x$ (cm), $x>1$.\\
	a) Nếu cạnh tăng thêm $2\text{ cm}$ thì thể tích khối lập phương mới là bao nhiêu? Viết kết quả dưới dạng đa thức và tính thể tích tăng thêm.\\
	b) Nếu cạnh giảm đi $1\text{ cm}$ thì thể tích giảm đi bao nhiêu?\\
	c) Áp dụng với $x=5$, tính thể tích tăng thêm và thể tích giảm đi trong hai trường hợp trên.
	\loigiai{
		a) Thể tích khối lập phương mới:
		\[ (x+2)^3=x^3+6x^2+12x+8\ \left(\text{cm}^3\right). \]
		Thể tích tăng thêm: $\left(x^3+6x^2+12x+8\right)-x^3=6x^2+12x+8\ \left(\text{cm}^3\right)$.\\
		b) Thể tích khi cạnh giảm $1\text{ cm}$: $(x-1)^3=x^3-3x^2+3x-1\ \left(\text{cm}^3\right)$.\\
		Thể tích giảm đi: $x^3-\left(x^3-3x^2+3x-1\right)=3x^2-3x+1\ \left(\text{cm}^3\right)$.\\
		c) Với $x=5$:\\
		Thể tích tăng thêm: $6\cdot 25+12\cdot 5+8=150+60+8=218\ \left(\text{cm}^3\right)$.\\
		Thể tích giảm đi: $3\cdot 25-3\cdot 5+1=75-15+1=61\ \left(\text{cm}^3\right)$.
	}
\end{bt}

\begin{bt}%[8D2V7-2]
	a) Viết biểu thức $-x^3+6x^2-12x+8$ dưới dạng lập phương của một hiệu.\\
	b) Viết biểu thức $x^6-3x^4y+3x^2y^2-y^3$ dưới dạng lập phương của một hiệu.\\
	c) Tính giá trị của biểu thức $-x^3+6x^2-12x+8$ tại $x=-8$.
	\loigiai{
		a) $-x^3+6x^2-12x+8=8-12x+6x^2-x^3=2^3-3\cdot 2^2\cdot x+3\cdot 2\cdot x^2-x^3=(2-x)^3$.\\
		b) $x^6-3x^4y+3x^2y^2-y^3=\left(x^2\right)^3-3\cdot\left(x^2\right)^2\cdot y+3\cdot x^2\cdot y^2-y^3=\left(x^2-y\right)^3$.\\
		c) Vì $-x^3+6x^2-12x+8=(2-x)^3$ nên tại $x=-8$ biểu thức có giá trị $(2-(-8))^3=10^3=1\,000$.
	}
\end{bt}

\begin{bt}%[8D2V7-3]
	a) Cho $x+y=5$ và $xy=6$. Tính giá trị của biểu thức $x^3+y^3$.\\
	b) Cho $x-y=3$ và $xy=10$. Tính giá trị của biểu thức $x^3-y^3$.
	\loigiai{
		a) Ta có $(x+y)^3=x^3+3x^2y+3xy^2+y^3=x^3+y^3+3xy(x+y)$.\\
		Suy ra $x^3+y^3=(x+y)^3-3xy(x+y)=5^3-3\cdot 6\cdot 5=125-90=35$.\\
		b) Ta có $(x-y)^3=x^3-3x^2y+3xy^2-y^3=x^3-y^3-3xy(x-y)$.\\
		Suy ra $x^3-y^3=(x-y)^3+3xy(x-y)=3^3+3\cdot 10\cdot 3=27+90=117$.
	}
\end{bt}
